% =====================================================================
%  preamble/commands.tex
%  Comandos y abreviaturas propias del documento.
% =====================================================================

% --- Subsecciones frecuentes ---
\newcommand{\AnalisisDescriptivo}{\subsubsection{Análisis descriptivo}}
\newcommand{\TablaDeDatos}{\subsubsection{Tabla de datos}}

% --- Resaltados ---
% \PalabraClave es SOLO para palabras clave: la palabra en negrita seguida de
% su valor, dentro de la línea de texto. No sirve como encabezado de bloque;
% para eso están \BloqueCorrido y \BloqueTitulo.
\newcommand{\PalabraClave}[1]{\medskip\textbf{#1}}
\newcommand{\Variable}[1]{\textbf{\texttt{#1}}}           % nombres de variables (monoespaciada)

% --- Encabezados de bloque ---
% Corrido: negrita + dos puntos + el texto, todo en el mismo párrafo.
%   \BloqueCorrido{Cómo se toman los datos}{Primero estandarizamos...}
\newcommand{\BloqueCorrido}[2]{\par\medskip\noindent\textbf{#1:}\space#2\par}
% Suelto: para cuando lo que sigue es una lista.
%   \BloqueTitulo{Análisis de las campanas}
%   \begin{itemize} ... \end{itemize}
\newcommand{\BloqueTitulo}[1]{\par\medskip\noindent\textbf{#1:}\par}
% \Interpretacion es un caso particular de \BloqueCorrido.
\newcommand{\Interpretacion}[1]{\BloqueCorrido{Interpretación}{#1}}
% \Respuesta no lleva argumento: lo que sigue (casi siempre una ecuación en
% display) va por fuera del macro, igual que con \BloqueTitulo.
\newcommand{\Respuesta}{\BloqueTitulo{Respuesta}}

% --- Título para gráficos ---
\newcommand{\TituloGrafico}[1]{\subsubsection{Gráfico: #1}}

% --- Encabezado y pie de página ---
\setlength{\headheight}{14pt}   % evita el warning de fancyhdr
\pagestyle{fancy}
\fancyhf{}
\fancyhead[L]{\small\textcolor{gris}{Trabajo Práctico de Estadística}}
\fancyhead[R]{\small\textcolor{gris}{Quiñonez -- Di Matteo}}
\fancyfoot[C]{\small\thepage}
\renewcommand{\headrulewidth}{0.4pt}

% --- Espaciado de secciones ---
\titlespacing*{\section}{0pt}{1.2em}{0.6em}
